\section{Conclusion}
\label{sec:conclusion}

Threshold cryptography has been ready for consumer deployment for years.
What has been missing is the product design that makes it invisible.
Threshold UX closes this gap by mapping threshold primitives to product
patterns that users already understand: approval flows, recovery contacts,
shared workspaces, spending limits, and compliance gates.

The key contributions of this paper are:

\begin{enumerate}
    \item \textbf{Threshold as UX primitive.} A formal mapping from
    threshold cryptographic operations (FROST signatures, Shamir sharing,
    threshold decryption, TFHE policy evaluation) to product interaction
    patterns (approvals, recovery, workspaces, treasury, compliance).

    \item \textbf{Social recovery without seed phrases.} A guardian-based
    recovery protocol that replaces 24-word mnemonics with ``3 of 5 trusted
    contacts,'' including guardian rotation via proactive secret sharing.

    \item \textbf{Progressive security.} A tiered security model that
    starts at 1-of-1 and upgrades to threshold control as value increases,
    with downgrade resistance preventing unilateral weakening.

    \item \textbf{UX interaction patterns.} Concrete design patterns
    (notification-driven approval, progress indicators, delegation,
    timeout escalation) that hide cryptographic ceremony behind familiar
    product metaphors.

    \item \textbf{Security equivalence proof.} A demonstration that the
    UX layer preserves the security guarantees of the underlying threshold
    primitives, introducing no additional attack surface beyond social
    engineering, which is bounded by guardian diversity.
\end{enumerate}

Threshold cryptography is not hard to use. It has been hard to present.
That problem is now solved.

\begin{thebibliography}{99}

\bibitem{shamir1979share}
A.~Shamir, ``How to share a secret,'' \emph{Communications of the ACM}, vol.~22, no.~11, pp.~612--613, 1979.

\bibitem{gennaro2018fast}
R.~Gennaro and S.~Goldfeder, ``Fast multiparty threshold ECDSA with fast trustless setup,'' in \emph{ACM CCS}, 2018.

\bibitem{komlo2020frost}
C.~Komlo and I.~Goldberg, ``FROST: Flexible round-optimized Schnorr threshold signatures,'' in \emph{SAC}, Springer, 2020.

\bibitem{canetti2021cggmp}
R.~Canetti, R.~Gennaro, S.~Goldfeder, N.~Makriyannis, and U.~Peled, ``UC non-interactive, proactive, threshold ECDSA with identifiable aborts,'' in \emph{ACM CCS}, 2020.

\bibitem{gentry2009fhe}
C.~Gentry, ``Fully homomorphic encryption using ideal lattices,'' in \emph{STOC}, ACM, 2009.

\bibitem{chillotti2020tfhe}
I.~Chillotti, N.~Gama, M.~Georgieva, and M.~Izabach\`ene, ``TFHE: Fast fully homomorphic encryption over the torus,'' \emph{Journal of Cryptology}, vol.~33, pp.~34--91, 2020.

\bibitem{herzberg1995proactive}
A.~Herzberg, S.~Jarecki, H.~Krawczyk, and M.~Yung, ``Proactive secret sharing or: How to cope with perpetual leakage,'' in \emph{CRYPTO}, Springer, 1995.

\bibitem{shapiro2011crdt}
M.~Shapiro, N.~Pregui\c{c}a, C.~Baquero, and M.~Zawirski, ``Conflict-free replicated data types,'' in \emph{SSS}, Springer, 2011.

\bibitem{blaze1998divertible}
M.~Blaze, G.~Bleumer, and M.~Strauss, ``Divertible protocols and atomic proxy cryptography,'' in \emph{EUROCRYPT}, Springer, 1998.

\bibitem{bonneau2012quest}
J.~Bonneau, C.~Herley, P.~C.~van~Oorschot, and F.~Stajano, ``The quest to replace passwords,'' in \emph{IEEE S\&P}, 2012.

\bibitem{eskandari2018first}
S.~Eskandari, D.~Barrera, E.~Stobert, and J.~Clark, ``A first look at the usability of Bitcoin key management,'' in \emph{NDSS Workshop on Usable Security}, 2018.

\bibitem{shoup2000practical}
V.~Shoup, ``Practical threshold signatures,'' in \emph{EUROCRYPT}, Springer, 2000.

\end{thebibliography}

